\documentclass[10pt,a4paper]{amsart}
\usepackage[margin=19mm]{geometry}
\usepackage{amsmath,amssymb,mathtools}
\usepackage[scr=rsfs,cal=euler]{mathalfa}
\usepackage{xfrac}
\usepackage[unicode,hidelinks,bookmarks=false]{hyperref}
\newcommand{\ie}{i.e.\ }
\newcommand{\cf}{cf.\ }
\newcommand{\resp}{respectively\ }
\newcommand{\Subsection}{\subsection}
\providecommand{\nobreakdash}{\nobreak\hbox{-}}
\setlength{\emergencystretch}{2em}
\multlinegap0pt
\usepackage{bm}
\usepackage{bbm}
\usepackage{dsfont}
\usepackage{xspace}
\usepackage{cancel}
\usepackage{mathtools}
\usepackage{amsthm}
\makeatletter
\@ifundefined{theorem}{\newtheorem{theorem}{Theorem}}{}
\makeatother
\title[A solution to Haagerup's problem and positive Hahn-Banach se]{A solution to Haagerup's problem and positive Hahn-Banach separation theorems in operator algebras}
\author{Ikhan Choi}
\date{Source: arXiv:2501.16832v3 (CC BY 4.0)}
\begin{document}
\maketitle

\noindent\emph{Source and licence.} A solution to Haagerup's problem and positive Hahn-Banach separation theorems in operator algebras. Version arXiv:2501.16832v3 (CC BY 4.0), distributed under the Creative Commons Attribution 4.0 International licence, as recorded in the arXiv licence metadata for this exact version and shown on the abstract page https://arxiv.org/abs/2501.16832v3. Full rights notice: licenses/arxiv-2501.16832-CC-BY-4.0.txt.\par\smallskip

\noindent\textit{Editorial scope.} Counted unit: the Theorem whose source label is \texttt{None} in the article (source0.tex lines 285--289, source label \texttt{None}), delivered together with the article's own complete original proof of it (source0.tex lines 290--337, one \texttt{proof} environment of 48 lines). No mathematics is written, repaired or re-derived: statement and proof are the article's own text line for line. Required context retained uncounted: none. Every definition, display and label of those lines is retained unchanged. Omitted from this package and not redistributed: 1--284, 338--389. Nothing inside a retained excerpt is omitted or altered: every retained body is delivered exactly as the article's lines are written, so no omission receipt of any kind accompanies this package. Citations: the retained excerpts carry no citation command at all, so no bibliography entry of the article is transcribed and no reference is redistributed. Reference adaptation: none was needed and none was made; every reference inside the delivered unit resolves inside it, so no unresolved reference or citation is produced. Printed slips kept literal: no printed word, symbol, hypothesis or number of the counted statement or of its proof was altered. Layout and class: the article's own class is \texttt{amsart} with packages including geometry, mathdesign, hyperref, setspace; the wrapper is the lane's pinned \texttt{amsart} editorial wrapper at 10pt on A4, which restates the theorem-like environments the retained text needs and adds the article's own macro definitions that the retained text needs (none), with extra wrapper packages (bm, bbm, dsfont, xspace, cancel, mathtools). The delivered numbering is the unit's own. Assets: no retained span contains a figure environment, a graphics-inclusion command, a PDF-page inclusion, a table or a TikZ picture; no image file is copied into this package, the package declares no asset dependency and no image file is redistributed with it. Licence: version arXiv:2501.16832v3 of this article is distributed under the Creative Commons Attribution 4.0 International licence, as recorded in the arXiv licence metadata for this exact version and shown on the abstract page https://arxiv.org/abs/2501.16832v3. A full rights notice is delivered at licenses/arxiv-2501.16832-CC-BY-4.0.txt. Characters: every retained excerpt is pure ASCII, so the delivered engine, XeLaTeX with its default Latin Modern text font, reports no missing character.
\begin{theorem}
Let $A$ be a C$^*$-algebra, and consider the dual pair $(A^{*sa},A^{sa})$.
If $F^*$ is a weakly$^*$ closed convex hereditary subset of $A^{*+}$, then $F^*=(F^*)^{r+r+}$.
Equivalently, if $\omega'\in A^{*+}\setminus F^*$, then there is $a\in A^+$ such that $\omega'(a)>1$ and $\omega(a)\le1$ for $\omega\in F^*$.
\end{theorem}

\begin{proof}
As above, we will prove $(\overline{F^*-A^{*+}})^+\subset F^*$.
Define
\[G^*:=\left\{\omega\in A^{*sa}:\begin{tabular}{c}
there is $\psi\in A^{*+}$, and there is $\varphi_\delta\in F^*$ \\
for any sufficiently small $\delta>0$,
such that\\
$\|\psi\|\le1$, $\|\varphi_\delta\|\le\delta^{-1}$, and $\omega\le\varphi_\delta+\delta^{\frac12}\psi$
\end{tabular}\right\}.\]
It suffices to show $F^*-A^{*+}\subset G^*$, $G^{*+}\subset F^*$, and $\overline{G^*}\subset G^*$.

Let $\omega\in F^*-A^{*+}$.
Take $\varphi\in F^*$ such that $\omega\le\varphi$.
Define, for $\delta>0$,
\[\psi:=\frac{[\omega]}{1+\|\omega\|}+\frac\varphi{(1+\|\omega\|)(1+\|\varphi\|)},\qquad\varphi_\delta:=\theta(f_\delta(\theta^{-1}(\varphi))),\]
where $\theta$ is the commutant Radon-Nikodym map associated to $\psi$.
The norm conditions $\|\psi\|\le1$ and $\|\varphi_\delta\|\le\delta^{-1}$ are easily checked.
For sufficiently small $\delta>0$ such that $\|\theta^{-1}(\omega)\|\le1+\|\omega\|\le2^{-1}\delta^{-\frac14}$ and $\delta\le1$, we have
\[\theta^{-1}(\omega)\le f_\delta(\theta^{-1}(\omega))+\delta^{\frac12}\le f_\delta(\theta^{-1}(\varphi))+\delta^{\frac12},\]
so $\omega\le\varphi_\delta+\delta^{\frac12}\psi$ and $\omega\in G^*$.

Let $\omega\in G^{*+}$.
Take $\psi\in A^{*+}$ and $\varphi_\delta\in F^*$ such that $\|\psi\|\le1$, $\|\varphi_\delta\|\le\delta^{-1}$, and $\omega\le\varphi_\delta+\delta^{\frac12}\psi$, for any sufficiently small $\delta>0$.
Let $\psi_\delta:=\omega+\delta\varphi+\psi$, and let $\theta_\delta$ be the associated commutant Radon-Nikodym map.
For any fixed $\delta'>0$, since $0\le\theta_\delta^{-1}(\omega)\le1$, we have
\begin{align*}
0&\le(1+\delta')^{-1}\theta_\delta^{-1}(\omega)\le f_{\delta'}(\theta_\delta^{-1}(\omega))\le f_{\delta'}(\theta_\delta^{-1}(\varphi_\delta+\delta^{\frac12}\psi))\\
&\le f_{\delta'}(\theta_\delta^{-1}(\varphi_\delta)+\delta^{\frac12})\le f_{\delta'}(\theta_\delta^{-1}(\varphi_\delta))+\delta^{\frac12},
\end{align*}
and it implies $0\le(1+\delta')^{-1}\omega\le\theta_\delta(f_{\delta'}(\theta_\delta^{-1}(\varphi_\delta)))+\delta^{\frac12}\psi_\delta$.
Since $\|\psi_\delta\|\le\|\omega\|+2$ is bounded and $\theta_\delta(f_{\delta'}(\theta_\delta^{-1}(\varphi_\delta)))\in F^*$ is also bounded for fixed $\delta'$, by considering the limit along a cofinal ultrafilter on the set of $\delta$, we have $(1+\delta')^{-1}\omega\in F^*$, so $\delta'\to0$ gives $\omega\in F^*$.

To show $G^*$ is weakly$^*$ closed, we claim for any $r>0$ that 
\[\overline{(F^*-A^{*+})\cap A_{2r}^*}\subset G^*,\qquad G^*\cap A_r^*\subset\overline{(F^*-A^{*+})\cap A_{2r}^*},\]
where $A_r^*:=\{\omega\in A^*:\|\omega\|\le r\}$.
If these are true, then
\[G^*\cap A_r^*=\overline{(F^*-A^{*+})\cap A_{2r}^*}\cap A_r^*\]
is weakly$^*$ closed and convex in $A^*$ for all $r>0$, so the Krein-\v Smulian theorem shows that $G^*$ is weakly$^*$ closed.

Let $\omega_i\in(F^*-A^{*+})\cap A_{2r}^*$ be a net such that $\omega_i\to\omega$ weakly$^*$ in $A^*$.
Following the proof of $F^*-A^{*+}\subset G^*$, we can take in the same way $\psi_i\in A^{*+}$ and $\varphi_{i,\delta}\in F^*$ such that $\|\psi_i\|\le1$, $\|\varphi_{i,\delta}\|\le\delta^{-1}$, and $\omega_i\le\varphi_{i,\delta}+\delta^{\frac12}\psi_i$, for uniformly sufficiently small $\delta$ such that $1+2r\le2^{-1}\delta^{-\frac14}$ because $\|\omega_i\|$ is bounded by $2r$.
Since the three conditions are preserved by the weak$^*$ convergence, taking the limit along a cofinal ultrafilter on the index set of $i$, we can obtain limit points $\psi$ and $\varphi_\delta$ so that $\omega\in G^*$.

Let $\omega\in G^*\cap A_r^*$.
Take $\psi\in A^{*+}$ and $\varphi_\delta\in F^*$ with $\|\psi\|\le1$, $\|\varphi_\delta\|\le\delta^{-1}$, and $\omega\le\varphi_\delta+\delta^{\frac12}\psi$, for any sufficiently small $\delta>0$.
If $\delta^{\frac12}<r$, then $\omega-\delta^{\frac12}\psi\in(F^*-A^{*+})\cap A_{2r}^*$ converges to $\omega$ weakly$^*$ in $A^*$ as $\delta\to0$, we have $\omega\in\overline{(F^*-A^{*+})\cap A_{2r}^*}$.
This completes the proof.
\end{proof}

\end{document}
